% =============================================================================
%  text/HISTORICAL BACKGROUND/background.tex
%  Demonstrates narrative paragraphs and a \citet citation.
%  Content is generic placeholder text -- replace it entirely.
%  演示叙述段落与 \citet 引用；内容为通用占位，请整体替换。
% =============================================================================

This subsection demonstrates how to write narrative paragraphs and cite
sources. All sentences below are placeholder text and bear no relation to
any research topic.

Background sections typically begin with a broad statement, then narrow down
to the specific context of the paper. Sources are cited as needed. For a
textual citation, use \texttt{\textbackslash citet}, as in
\citet{author2018book}; for a parenthetical citation, use
\texttt{\textbackslash citep}, as in \citep{author2020example}.

% 中文段落示例（可选）
% 背景部分通常先作总体陈述，再聚焦到论文的具体语境，并在需要处引用
% 文献。请以你自己的内容替换本段。
Chinese-language sources can be cited in the same way \citep{zhang2018example},
as long as the \texttt{.bib} file is saved in UTF-8.

This is the end of the placeholder background subsection.
